\section{Complete Weyl-basis marginals and LU rigidity}\label{sec:rigidity}

We first isolate the linear-algebra statement that recovers a diagonal
tensor's coordinate axes.  It is the complete-Weyl-basis counterpart of the
three-Pauli argument in \cite{Rains1999,VanDenNestDehaeneDeMoor2005}.
The statement itself is classical: it is Jennrich's uniqueness theorem for
trilinear decompositions, published in \cite[\S V]{Harshman1970} and there
credited to Jennrich, whose conclusion that the factor matrices agree up to
one permutation and compensating diagonal scalings is the monomial property
below.  It is the full-column-rank case of Kruskal's condition
\cite{Kruskal1977}, extended to arbitrarily many factors in
\cite{SidiropoulosBro2000}; see \cite[\S3.2]{KoldaBader2009} for the
attribution and for a survey of the uniqueness theory.  We give a short
self-contained proof because the contraction argument is the one used in
Proposition~\ref{prop:full-weyl-marginal}, and claim no novelty for either
the statement or the proof.

\begin{lemma}[Diagonal-tensor axes]\label{lem:diagonal-axes}
Let \(r\geq3\), let \(E_1,\ldots,E_r\) be \(N\)-dimensional complex inner
product spaces, and choose orthonormal bases
\(\{e_{ij}:1\leq j\leq N\}\).  If
\[
 T=\sum_{j=1}^N\lambda_j
 e_{1j}\otimes\cdots\otimes e_{rj},\qquad \lambda_j\ne0,       \tag{3.1}
\]
then the coordinate axes in every \(E_i\) are intrinsic to \(T\).
Consequently, a product of invertible linear maps carrying one tensor of
the form \((3.1)\) to another is monomial in the displayed bases.
\end{lemma}

\begin{proof}
For nonzero \(x=\sum_jx_je_{1j}^*\), contract the first factor against
\(x\).  The result is
\[
 T_x=\sum_j\lambda_jx_j e_{2j}\otimes\cdots\otimes e_{rj}.
\]
Flattening from \(E_2\) to \(E_3\otimes\cdots\otimes E_r\) gives a matrix
of rank \(\#\{j:x_j\ne0\}\).  Thus \(T_x\) has rank one exactly when
the ray of \(x\) is a dual coordinate axis.  Product equivalence preserves
this projective rank-one contraction locus.  The inner products identify
the recovered dual axes with the coordinate axes in the first factor.
Apply the same argument in every factor.
\end{proof}

\AMEFigureTwo

To avoid repeatedly stating the matched-basis condition, call an operator on
\(r\geq3\) qudits \emph{Weyl-diagonal with complete local support} if there
is a \(q^2\)-element set \(P\), a distinguished element \(0\in P\), and
bijections \(f_i:P\to\F_q^2\) with \(f_i(0)=0\), for which
\[
 R=\sum_{v\in P}\lambda_v
       \bigotimes_{i=1}^r W_{f_i(v)},\qquad \lambda_v\ne0.       \tag{3.2}
\]
The coefficients need not be equal.  What matters is that the same
\(q^2\)-element index set runs bijectively through the complete local
Weyl basis at every party.  Neither \(P\) nor the bijections need be
\(\F_q\)-linear.

\begin{proposition}[Complete Weyl-basis marginal criterion]
\label{prop:full-weyl-marginal}
Let \(R\) and \(R'\) be Weyl-diagonal operators with complete local support
on the same
\(r\geq3\) parties, possibly with different indexing sets, bijections,
and coefficients.  If
\[
 R'=(U_1\otimes\cdots\otimes U_r)R
       (U_1\otimes\cdots\otimes U_r)^\dagger ,
\]
then every \(U_i\) is Clifford.
\end{proposition}

\begin{proof}
View \(R\) and \(R'\) as tensors in the local Hilbert--Schmidt spaces.
After replacing \(W_v\) by the orthonormal basis \(q^{-1/2}W_v\) and
absorbing the common normalization into the coefficients, both have the
form in Lemma~\ref{lem:diagonal-axes}, with \(N=q^2\).
Local conjugation acts invertibly on each factor, so it must permute the
Weyl axes.  It fixes the identity operator and hence the identity axis.
It therefore permutes the nonidentity projective Weyl operators, which
is exactly the single-qudit Clifford condition.
\end{proof}

\begin{corollary}[Marginal-cover rigidity]
\label{cor:full-weyl-cover}
Let \(\lvert\psi\rangle,\lvert\phi\rangle\) be \(n\)-qudit states and
suppose a product unitary carries \(\lvert\psi\rangle\) to
\(\lvert\phi\rangle\) up to phase, any party permutation having first been
absorbed by relabelling \(\lvert\phi\rangle\), so that the two states are
indexed by the same parties.  If every
party belongs to a set \(S\) of at least three parties for which both
reduced operators \(\rho^\psi_S\) and \(\rho^\phi_S\) are Weyl-diagonal
with complete local support, then every local factor is Clifford.
\end{corollary}

\begin{proof}
Partial trace gives a product-unitary equivalence of the two reduced
operators on each chosen \(S\).  Apply
Proposition~\ref{prop:full-weyl-marginal}; the chosen sets cover all
parties.
\end{proof}

\begin{proposition}[Stabilizer-AME support theorem]
\label{prop:stabilizer-ame-support}
Let \(\lvert\psi\rangle\) be a stabilizer
\(\AME(2m,q)\) state, and let \(A\) be any \(m+1\) parties.  Then
\[
 |L(A)|=q^2.
\]
Every nonidentity label in \(L(A)\) has support exactly \(A\), and for
each \(i\in A\) the coordinate projection
\[
 p_i:L(A)\longrightarrow\F_q^2
\]
is bijective.
\end{proposition}

\begin{proof}
The projective label group \(L\) has order \(q^{2m}\).  Restrict its
labels to the \(m-1\) parties outside \(A\).  The target has order
\(q^{2(m-1)}\), and the kernel is \(L(A)\); hence
\[
 |L(A)|\geq q^{2m}/q^{2(m-1)}=q^2.                    \tag{3.3}
\]
Fix \(i\in A\).  A label in the kernel of \(p_i\) is supported on the
\(m\)-set \(A\setminus\{i\}\).  The AME condition and \((2.3)\) exclude
every such nonidentity stabilizer, so \(p_i\) is injective.  Its codomain
has order \(q^2\), giving the reverse inequality in \((3.3)\) and hence
bijectivity.  The same no-small-support argument shows that every
nonidentity element of \(L(A)\) has full support \(A\).
\end{proof}

\begin{proposition}[A stabilizer-AME marginal fixes the local Weyl axes]
\label{prop:marginal-axes}
Let \(\lvert\psi\rangle,\lvert\phi\rangle\) be stabilizer
\(\AME(2m,q)\) states.  If a product unitary on an \((m+1)\)-set \(A\)
carries \(\rho^\psi_A\) to \(\rho^\phi_A\), then every local factor on
\(A\) is Clifford.
\end{proposition}

\begin{proof}
By \((2.3)\),
\[
 \rho^\psi_A=q^{-|A|}
   \sum_{\ell\in L_\psi(A)}\widetilde W_\ell .
\]
Index this sum by \(P=L_\psi(A)\).  Proposition
\ref{prop:stabilizer-ame-support} says that every coordinate projection
\(P\to\F_q^2\) is bijective.  Relative to the fixed Weyl convention, each
stabilizer lift contributes a nonzero coefficient.  Thus
\(\rho^\psi_A\) is Weyl-diagonal with complete local support in the sense of
\((3.2)\), and the same holds for \(\rho^\phi_A\).  Apply
Proposition~\ref{prop:full-weyl-marginal}.
\end{proof}

\begin{proof}[Proof of Theorem~\ref{thm:lu-lc-rigidity}]
Absorb the party permutation by relabelling \(\lvert\phi\rangle\).  Partial
trace gives a product-unitary equivalence
\(\rho^\psi_A\sim\rho^\phi_A\) for every \((m+1)\)-set \(A\).  Given a
party \(i\), choose such a set containing it and apply
Proposition~\ref{prop:marginal-axes}.  Thus every \(U_i\) is Clifford.
\end{proof}

\section{Minimum-support transition maps and encoder consequences}
\label{sec:atlas}

The same support calculation recovers standard quantum-MDS structure.
If \(|A|\leq m\), the AME condition gives \(L(A)=\{0\}\).  If
\(|A|>m\), then \(\rho^\psi_{A^c}\) is maximally mixed, so purity of the
global state gives
\(\Tr((\rho^\psi_A)^2)=q^{-(2m-|A|)}\).  On the other hand, \((2.3)\)
and Weyl orthogonality give
\(\Tr((\rho^\psi_A)^2)=q^{-|A|}|L(A)|\).  Hence, for every party set
\(A\),
\[
 |L(A)|=q^{2\max(|A|-m,0)}.                              \tag{4.1}
\]
M\"obius inversion is the stabilizer-state specialization of the
quantum-MDS Shor--Laflamme distribution of Huber and Grassl
\cite[Theorems~8 and 9]{HuberGrassl2020}; in particular the number of
weight-\(m+1\) projective stabilizers is
\(\binom{2m}{m+1}(q^2-1)\).  The following systematic form strengthens
minimum-weight generation: recognition needs only \(m\) of these supports,
and that number is optimal among spanning families of minimum-support
subgroups.

\begin{proposition}[Half-set systematic decomposition]
\label{prop:half-set-direct-sum}
Fix an \(m\)-party set \(B\), and let \(P_j=\F_q^2\) be the local label
space at party \(j\).  The coordinate projection
\[
 \operatorname{pr}_{B^c}:L\longrightarrow
 \bigoplus_{j\notin B}P_j
\]
is an isomorphism of additive groups.  Hence there is a unique
\(\F_p\)-linear map
\[
 K_B:\bigoplus_{j\notin B}P_j\longrightarrow
     \bigoplus_{i\in B}P_i
\]
whose graph is \(L\).  For the direct-sum trace-symplectic forms it
satisfies
\[
 [K_Bx,K_By]_B=-[x,y]_{B^c}.
\]
After local \(\F_p\)-bases are chosen, row reduction therefore puts any
row-generator check matrix for \(L\), ordered as \(B,B^c\), in systematic
form \([K_B^{\mathsf T}\mid I]\); here \(K_B\) acts on column vectors.
Moreover, if \(I\subseteq B\), \(J\subseteq B^c\), and
\(|I|=|J|\), then the party-block submatrix
\[
 K_{I,J}:\bigoplus_{j\in J}P_j\longrightarrow
          \bigoplus_{i\in I}P_i
\]
is an isomorphism.
Conversely, a maximal isotropic subspace of the form
\(\operatorname{graph}(K_B)\) is the label space of a stabilizer AME state
exactly when all these party-aligned square block maps are isomorphisms.
Moreover, \(K_{I,J}\) is invertible if and only if the complementary block
\(K_{B\setminus I,B^c\setminus J}\) is invertible.  Thus it suffices to test
one representative of each complementary pair.

For each \(j\notin B\), put
\(A_j=B\cup\{j\}\).  The inverse image of the \(j\)-th coordinate
summand is \(L(A_j)\), and therefore
\[
 L=\bigoplus_{j\notin B}L(A_j).
\]
Thus these \(m\) minimum-support subgroups determine the full stabilizer
label group.  Moreover, no collection of fewer than \(m\)
minimum-support subgroups can span \(L\).
\end{proposition}

\begin{proof}
The kernel of \(\operatorname{pr}_{B^c}\) is \(L(B)=\{0\}\) by the AME
support condition.  Its domain and codomain both have order \(q^{2m}\),
so it is an isomorphism.  This gives the graph map \(K_B\) and the
systematic check-matrix form.  Since \(L\) is isotropic, evaluating the
ambient direct-sum symplectic form on \((K_Bx,x)\) and \((K_By,y)\) gives
the displayed identity.

For the block statement, let
\(S=I\cup(B^c\setminus J)\), so \(|S|=m\).  Projection \(L\to
\bigoplus_{s\in S}P_s\) is an isomorphism by the same kernel-and-order
argument.  If \(x\) is supported on \(J\) and \(K_{I,J}x=0\), then
\((K_Bx,x)\in L\) projects to zero on \(S\).  Thus \(x=0\).
The source and target of \(K_{I,J}\) have equal dimension, so the block
map is an isomorphism.

Conversely, every \(m\)-party set has the form
\(I\cup(B^c\setminus J)\) with \(|I|=|J|\).  The kernel of projection
from \(\operatorname{graph}(K_B)\) to this set is naturally
\(\ker K_{I,J}\).  If all the displayed block maps are isomorphisms, every
half-set projection is therefore injective and hence bijective.  Equivalently,
there is no nonzero stabilizer label supported on at most \(m\) parties,
which is the stabilizer AME condition.

For complementarity, \(K_{I,J}\) is invertible exactly when
\(K_B(\bigoplus_{j\in J}P_j)\) is transverse to
\(\bigoplus_{i\in B\setminus I}P_i\).  Taking symplectic orthogonal
complements and using the anti-symplectic identity for \(K_B\) gives the
same condition for \(K_{B\setminus I,B^c\setminus J}\).

The inverse image of the \(j\)-th coordinate
summand consists exactly of the labels supported on \(B\cup\{j\}\).
Pulling back the coordinate direct sum gives the displayed decomposition.

Every minimum-support subgroup has order \(q^2\) by
Proposition~\ref{prop:stabilizer-ame-support}.  A sum of \(r\) such
subgroups has order at most \(q^{2r}\), whereas \(|L|=q^{2m}\).
Spanning \(L\) therefore requires \(r\geq m\).
\end{proof}

We now prove the transition-map classification stated in the introduction.
Regard each local label group
\[
 P_i=\F_q^2
\]
as a \(2e\)-dimensional trace-symplectic space over \(\F_p\).  For every
\((m+1)\)-set \(A\) and \(i,j\in A\), the support theorem defines the
operator-pushing transition
\[
 M_A(i,j)=p_jp_i^{-1}:P_i\longrightarrow P_j.
\]
\AMEFigureOne

For two states \(\psi,\phi\) on the same labelled parties, a
\emph{symplectic equivalence of the transition systems} is a tuple
\(F_i\in\operatorname{Sp}(P_i)\) satisfying
\[
 F_jM_A^\psi(i,j)=M_A^\phi(i,j)F_i
 \quad\text{for every \(A\) and \(i,j\in A\)}.           \tag{4.2}
\]
Here
\(\operatorname{Sp}(P_i)=\operatorname{Sp}_{2e}(\F_p)\) is the action
induced by the standard
additive local Clifford group.

For one state, define
\[
\begin{split}
 \Gamma_\psi
  &=\{[U]\in\operatorname{PU}(q^{2m}):
       U=U_1\otimes\cdots\otimes U_{2m},\
       U\lvert\psi\rangle\in\mathbb C\lvert\psi\rangle\},\\
 \operatorname{Aut}_{\mathrm{tr}}(\psi)
  &=\{(F_i)\in\prod_i\operatorname{Sp}(P_i):
       F_jM_A(i,j)=M_A(i,j)F_i\ \text{for all }A,i,j\in A\}.
\end{split}
\]
Here the subscript \(\mathrm{tr}\) refers to the transition system; this
group is also the linear image of the fixed-party transversal symmetries.
Theorem~\ref{thm:lu-lc-rigidity} and the Clifford sequence \((2.1)\)
give a well-defined homomorphism
\[
 \lambda_\psi:\Gamma_\psi\longrightarrow
 \operatorname{Aut}_{\mathrm{tr}}(\psi),\qquad
 [U_1\otimes\cdots\otimes U_{2m}]\longmapsto(\overline U_i)_i,
\]
where \(\overline U_i\) is the induced symplectic label map.  A product
decomposition is unique up to one-site scalars, which do not change
\(\overline U_i\).  The injection
\(L_\psi\hookrightarrow\Gamma_\psi\) sends a label
\(\ell\) to the projective Pauli \([W_\ell]\).

\begin{lemma}[Pauli phase correction]\label{lem:pauli-phase-correction}
Let \(\psi\) and \(\phi\) be stabilizer states with Lagrangian label groups
\(L_\psi\) and \(L_\phi\).  If a product Clifford \(U\) induces a
symplectic label map carrying \(L_\psi\) onto \(L_\phi\), then there is a
product Pauli \(W_v\) such that
\(W_vU\lvert\psi\rangle\in\mathbb C\lvert\phi\rangle\).
If \(L_\phi\) is the label space of a stabilizer \(\AME(2m,q)\) state,
then for any prescribed \(m\)-party set \(C\), the correcting label
\(v\) may be chosen uniquely with support contained in \(C\).  No prescribed
set of fewer than \(m\) parties can support corrections for every possible
character discrepancy.
\end{lemma}

\begin{proof}
The state \(U\lvert\psi\rangle\) is stabilized on the same Pauli axes as
\(\lvert\phi\rangle\), possibly with a different stabilizer character.
The trace-symplectic pairing identifies the ambient label space modulo
\(L_\phi=L_\phi^\perp\) with the character dual of \(L_\phi\).  Choose
\(v\) whose pairing character is the inverse of the discrepancy.
Conjugation by the product Pauli \(W_v\) cancels it.  The corrected state
and \(\lvert\phi\rangle\) therefore lie in the same one-dimensional joint
stabilizer eigenspace.

For the final statement, projection \(L_\phi\to\bigoplus_{i\in C}P_i\)
is an isomorphism by Proposition~\ref{prop:half-set-direct-sum}.  The
trace-symplectic form on the latter space is nondegenerate.  Hence pairing
with labels supported on \(C\) identifies
\(\bigoplus_{i\in C}P_i\) with the full character dual of \(L_\phi\).
Exactly one such label realizes the inverse discrepancy.
If a prescribed support has \(r<m\) parties, its label space has only
\(q^{2r}\) elements, whereas the character group of \(L_\phi\) has
\(q^{2m}\).  It therefore cannot realize every discrepancy.
\end{proof}

\begin{proof}[Proof of Theorem~\ref{thm:atlas-classification}]
First fix the party labels.  A block symplectic map
\(F=\bigoplus_iF_i\) carrying \(L_\psi\) to \(L_\phi\) carries every
supported subgroup \(L_\psi(A)\) to \(L_\phi(A)\).  Writing these
subgroups as graphs through any projection \(p_i\) gives exactly
\((4.2)\).

Conversely, suppose \((4.2)\) holds.  Fix \(A\) and \(i\in A\).  Every
element of \(L_\psi(A)\) is uniquely determined by its \(i\)-coordinate,
and the equations with \(j\in A\) show that its image under \(F\) is the
unique element of \(L_\phi(A)\) with \(i\)-coordinate \(F_i p_i(\ell)\).
Thus \(F L_\psi(A)=L_\phi(A)\).  Since the \((m+1)\)-supported subgroups
generate the full label groups by
Proposition~\ref{prop:half-set-direct-sum}, \(F L_\psi=L_\phi\).

Sequence \((2.1)\) supplies local Clifford lifts of the \(F_i\), and
Lemma~\ref{lem:pauli-phase-correction} supplies the product-Pauli
correction.  Hence an equivalence of transition systems lifts to an LC
equivalence.
The reverse implication follows
from the induced label action, and Theorem~\ref{thm:lu-lc-rigidity}
replaces LC by LU.  A party permutation is absorbed by relabelling first.

For automorphisms, \(\lambda_\psi\) is surjective by the same lifting and
phase-correction argument.  Its kernel consists of projective product
Paulis by \((2.1)\).  Such a Pauli fixes the state ray exactly when its
label lies in \(L_\psi^\perp=L_\psi\).  Thus the displayed injection of
\(L_\psi\) identifies it with \(\ker\lambda_\psi\), proving exactness.

\end{proof}

\begin{corollary}[Finite recognition bound]
\label{cor:finite-recognition}
Let \(\psi\) and \(\phi\) be stabilizer \(\AME(2m,q)\) states,
\(q=p^e\) and \(m\geq2\), given by check matrices for their additive
label spaces.  With the party labels fixed, choose any \(m\)-party set
\(B\), and write their half-set systematic maps as
\(K_B^\psi,K_B^\phi\).  The states are LU equivalent if and only if
there are \(F_i\in\operatorname{Sp}_{2e}(\F_p)\) satisfying
\[
 \left(\bigoplus_{i\in B}F_i\right)K_B^\psi
 =K_B^\phi\left(\bigoplus_{j\notin B}F_j\right).       \tag{4.3}
\]
This can be decided by testing at most
\[
 |\operatorname{Sp}_{2e}(\F_p)|
 =p^{e^2}\prod_{r=1}^e(p^{2r}-1)
\]
base symplectic frames.  In the arithmetic-operation model over \(\F_p\),
each candidate requires polynomial work in the check-matrix size.  The
stabilizer phases need not be supplied for this decision:
Lemma~\ref{lem:pauli-phase-correction} guarantees a character repair after
the label spaces are matched.  If the two characters are supplied, then for
any chosen \(m\)-party set the repair is the unique Pauli label on that
carrier realizing their discrepancy, and it is found by one linear solve.
For prime \(q\),
the number of base candidates is
\(|\mathrm{SL}_2(q)|=q(q^2-1)\).
In particular, fixed-label recognition is polynomial in \(m\) for fixed
local dimension \(q\).
With classical cubic matrix arithmetic, a concrete bound is
\[
 O\!\left((me)^3+
 |\operatorname{Sp}_{2e}(\F_p)|\,m^2e^3\right)
\]
\(\F_p\)-operations.

If the party relabelling is unknown, exhaustive search introduces at most a
factor \((2m)!\).  For one state,
\[
 |\Gamma_\psi|
 =q^{2m}|\operatorname{Aut}_{\mathrm{tr}}(\psi)|
 \leq q^{2m}|\operatorname{Sp}_{2e}(\F_p)|.
\]
\end{corollary}

\begin{proof}
Mapping the two graph subspaces into one another is exactly the block
equation \((4.3)\).  By Theorem~\ref{thm:atlas-classification} and
Lemma~\ref{lem:pauli-phase-correction}, its locally symplectic solutions
are exactly the LU equivalences.

Write \(K^\psi_{ij}\) and \(K^\phi_{ij}\) for the party blocks.  They are
invertible by Proposition~\ref{prop:half-set-direct-sum}.  Choose
\(i_0\in B\) and \(j_0\notin B\).  Once \(F_{j_0}\) is fixed, the
\((i,j_0)\)-blocks of \((4.3)\) determine every \(F_i\), and the
\((i_0,j)\)-blocks then determine every remaining \(F_j\).  Enumerate
\(F_{j_0}\in\operatorname{Sp}_{2e}(\F_p)\), test that the propagated
frames are symplectic, and check all \(m^2\) block equations.  Forming the
systematic maps and carrying out these checks is finite-field linear
algebra polynomial in the input size and gives the displayed arithmetic
bound.  The \(2m-1\) propagation equations form a spanning tree in the
complete bipartite party graph; the remaining
\(m^2-(2m-1)=(m-1)^2\) checks are its fundamental four-cycles.
Enumerating party permutations gives the stated factor.  The group-order
bound follows from the exact sequence in
Theorem~\ref{thm:atlas-classification} and injectivity of evaluation at the
base frame.
\end{proof}

\begin{corollary}[Prime-field recognition by fundamental cycles]
\label{cor:prime-systematic-recognition}
Let \(\psi,\phi\) be stabilizer \(\AME(2m,q)\) states with \(m\geq2\),
assume \(q=p\), and fix a labelled split \(B\sqcup B^c\).  Write
\(G=K_B^\psi\), \(H=K_B^\phi\), and choose
\(b\in B\) and \(c\in B^c\).  For \(i\ne b\) and \(j\ne c\), define
\[
 Q^G_{ij}
 =G_{bj}G_{ij}^{-1}G_{ic}G_{bc}^{-1}\in\operatorname{GL}_2(q),
\]
and define \(Q^H_{ij}\) similarly.  Then \(\psi\) and \(\phi\) are
fixed-party LU equivalent if and only if
\[
 \det G_{ij}=\det H_{ij}\quad\text{for all \(i,j\)}
\]
and there is an \(A\in\mathrm{SL}_2(q)\) such that
\[
 A Q^G_{ij}=Q^H_{ij}A
 \quad\text{for all \(i\ne b,\ j\ne c\)}.
\]
Thus prime-field recognition reduces to homogeneous linear equations in
the four entries of \(A\), followed by the single quadratic condition
\(\det A=1\).  Equivalently, the block-determinant array together with the
simultaneous \(\mathrm{SL}_2(q)\)-conjugacy class of the ordered fundamental
holonomies is a complete fixed-party LU invariant.  The corresponding
symmetry algebra and its five prime-field types are identified intrinsically
in Section~\ref{sec:endomorphism-algebra}.
\end{corollary}

\begin{proof}
Suppose first that
\[
 A_iG_{ij}=H_{ij}D_j
\]
for local \(A_i,D_j\in\mathrm{SL}_2(q)\).  Taking determinants gives the
block-determinant equalities.  Multiplying around the displayed four-cycle
gives \(Q^H_{ij}=A_bQ^G_{ij}A_b^{-1}\), so \(A=A_b\) satisfies the
intertwining equations.

Conversely, start from such an \(A\) and set
\[
 D_j=H_{bj}^{-1}AG_{bj},\qquad
 A_i=H_{ic}D_cG_{ic}^{-1}.
\]
The determinant equalities put every \(D_j\) and \(A_i\) in
\(\mathrm{SL}_2(q)\).  The base row and base column equations hold by
construction.  Each remaining equation \(A_iG_{ij}=H_{ij}D_j\) is
equivalent, after substitution, to
\(AQ^G_{ij}=Q^H_{ij}A\).  Hence the systematic matrices are locally
symplectically equivalent, and Corollary~\ref{cor:finite-recognition}
lifts this to LU equivalence.  Taking \(G=H\) gives the centralizer
statement.
\end{proof}

\subsection{Marginal certification and stochastic conversion}
\label{subsec:marginal-certification}

The half-set direct sum has two consequences outside the equivalence
problem.  The first says which reductions of a stabilizer AME state
determine it among all density operators and gives the matching parent
Hamiltonian; the second removes the unitarity hypothesis from
Theorem~\ref{thm:lu-lc-rigidity}.  For a party set \(A\), write
\(\sigma_A\) for the reduction of a density operator \(\sigma\) to \(A\),
and \(\|X\|_1=\Tr|X|\).

\begin{proposition}[Marginal certification]
\label{prop:marginal-certification}
Let \(\lvert\psi\rangle\) be a stabilizer \(\AME(2m,q)\) state with
\(m\geq2\), fix an \(m\)-party set \(B\), and for \(j\notin B\) put
\(A_j=B\cup\{j\}\) and
\[
 \Pi_j=q^{m-1}\rho^\psi_{A_j}\otimes I_{A_j^c}.
\]
\begin{enumerate}
\item Each \(\Pi_j\) is the orthogonal projector onto the joint \(+1\)
 eigenspace of the stabilizer lifts with labels in \(L(A_j)\).  The
 \(\Pi_j\) commute, and
 \(\prod_{j\notin B}\Pi_j=\lvert\psi\rangle\langle\psi\rvert\).
\item A density operator \(\sigma\) on the \(2m\) parties with
 \(\sigma_{A_j}=\rho^\psi_{A_j}\) for every \(j\notin B\) equals
 \(\lvert\psi\rangle\langle\psi\rvert\).  More generally, the reductions
 of \(\psi\) to \((m+1)\)-party sets \(C_1,\ldots,C_k\) determine \(\psi\)
 among all density operators if and only if \(L(C_1),\ldots,L(C_k)\) span
 \(L\); in particular \(k\geq m\).
\item Every reduction of \(\psi\) to at most \(m\) parties is maximally
 mixed, so no family of such reductions distinguishes \(\psi\) from the
 maximally mixed state.
\item \(H=\sum_{j\notin B}(I-\Pi_j)\) is a sum of \(m\) commuting
 projectors, each acting on \(m+1\) parties, with kernel
 \(\mathbb C\lvert\psi\rangle\) and spectral gap exactly one.  For every
 density operator \(\sigma\),
 \[
  1-\langle\psi\rvert\sigma\lvert\psi\rangle
  \leq\Tr(H\sigma)
  \leq\sum_{j\notin B}\tfrac12\bigl\|\sigma_{A_j}-\rho^\psi_{A_j}\bigr\|_1 .
 \]
\item A Hamiltonian that is a sum of terms acting on at most \(m\) parties
 each and has \(\lvert\psi\rangle\) as a ground state is a scalar.
\end{enumerate}
\end{proposition}

\begin{proof}
(i) By \((2.3)\) and Proposition~\ref{prop:stabilizer-ame-support},
\[
 q^{m-1}\rho^\psi_{A_j}\otimes I_{A_j^c}
 =q^{-2}\sum_{\ell\in L(A_j)}\widetilde W_\ell ,
\]
and the lifts on the right form the subgroup of the stabilizer group with
labels in \(L(A_j)\), of order \(q^2\).  The normalized sum over a finite
abelian group of unitaries is the projector onto its joint \(+1\)
eigenspace, so \(\Pi_j\) is that projector, and the \(\Pi_j\) commute
because the stabilizer group is abelian.  By
Proposition~\ref{prop:half-set-direct-sum}, addition is a bijection
\(\bigoplus_{j\notin B}L(A_j)\to L\), and the lift is multiplicative on
the stabilizer group; hence
\[
 \prod_{j\notin B}\Pi_j
 =q^{-2m}\sum_{\ell\in L}\widetilde W_\ell
 =\lvert\psi\rangle\langle\psi\rvert
\]
by \((2.2)\).

(ii) Let \(G=L(C_1)+\cdots+L(C_k)\) and
\(P_i=q^{-2}\sum_{\ell\in L(C_i)}\widetilde W_\ell\), a projector by the
argument of (i).  If \(\sigma_{C_i}=\rho^\psi_{C_i}\), then Weyl
orthogonality and \(|L(C_i)|=q^2\) give
\(\Tr(\sigma P_i)=q^{m-1}\Tr\bigl((\rho^\psi_{C_i})^2\bigr)=1\); since
\(0\leq P_i\leq I\), this forces \(P_i\sigma=\sigma\).  If the
\(L(C_i)\) span \(L\), then \(\prod_iP_i\) is the projector onto the joint
\(+1\) eigenspace of the whole stabilizer group, which is
\(\lvert\psi\rangle\langle\psi\rvert\), so
\(\sigma=\lvert\psi\rangle\langle\psi\rvert\sigma\); taking adjoints,
\(\sigma
 =\lvert\psi\rangle\langle\psi\rvert\sigma\lvert\psi\rangle\langle\psi\rvert
 =\lvert\psi\rangle\langle\psi\rvert\).
If instead \(G\neq L\), put
\[
 \tau=q^{-2m}\sum_{\ell\in G}\widetilde W_\ell .
\]
This is \(|G|q^{-2m}\) times the projector onto the joint \(+1\) eigenspace
of the lifts of \(G\), whose rank \(q^{2m}/|G|\) exceeds one; so \(\tau\)
is a density operator of rank greater than one.  Tracing out \(C_i^c\)
annihilates every lift not supported in \(C_i\), and
\(G\cap L(C_i)=L(C_i)\), so
\[
 \tau_{C_i}=q^{-2m}\cdot q^{m-1}\sum_{\ell\in L(C_i)}\widetilde W_\ell
 =\rho^\psi_{C_i}
\]
by \((2.3)\).  Thus the reductions do not determine \(\psi\).  The first
sentence is the case \(C_i=A_j\), where
Proposition~\ref{prop:half-set-direct-sum} gives the spanning, and its
final sentence gives \(k\geq m\).

(iii) is \((2.3)\) with the AME condition.

(iv) Commuting projectors make \(H\) diagonalizable with integer
eigenvalues, and the eigenvalue \(0\) has eigenspace
\(\bigcap_j\operatorname{range}\Pi_j=\operatorname{range}\prod_j\Pi_j
 =\mathbb C\lvert\psi\rangle\) by (i).  Hence
\(H\geq I-\lvert\psi\rangle\langle\psi\rvert\), which is the first
inequality.  The gap is attained: a Pauli \(W_v\) whose pairing character
is nontrivial on exactly one summand \(L(A_j)\) exists because the
characters of \(L\) are realized by labels modulo \(L\), and
\(W_v\lvert\psi\rangle\) has \(H=1\).  For the second inequality,
\(\Tr(\sigma(I-\Pi_j))
 =\Tr\bigl((\rho^\psi_{A_j}-\sigma_{A_j})P_j\bigr)\)
with \(P_j=q^{m-1}\rho^\psi_{A_j}\), using
\(\Tr(\rho^\psi_{A_j}P_j)=1\).  For a traceless Hermitian \(X\) and
\(0\leq P\leq I\), \(\Tr(XP)\leq\Tr X_+=\tfrac12\|X\|_1\).  Sum over
\(j\).

(v) Let \(H'=\sum_kh_k\) with \(h_k\) acting on a set \(S_k\) of at most
\(m\) parties.  By (iii),
\(\langle\psi\rvert h_k\lvert\psi\rangle=q^{-2m}\Tr h_k\), so
\(\langle\psi\rvert H'\lvert\psi\rangle=q^{-2m}\Tr H'\) is the mean of
the eigenvalues of \(H'\).  If \(\psi\) is a ground state, the least
eigenvalue equals the mean, so all eigenvalues coincide.
\end{proof}

For qubit stabilizer states, Wu, Yang, Wang, Wen, Qin, and Gao show that
the reductions to the supports of any \(n\) generators determine the state
among all density operators \cite[Theorem~2]{WuYangWangWenQin2015}.  The
AME support geometry lowers the count from \(2m\) to \(m\), fixes every
support to \(m+1\) parties, and shows that neither number can be reduced.
Generic pure states on \(n\) parties are determined by reductions to about
\(n/2+1\) parties \cite{JonesLinden2005}; here the locality is the same
but the family is explicit.  The certificate in (iv) is the stabilizer
witness of \cite{TothGuhne2005} applied to a frustration-free Hamiltonian,
and the state \(\sigma\) need not be pure, stabilizer, or AME, which
separates it from the approximate product-symmetry statements of
Section~\ref{sec:quantitative}.

\begin{corollary}[Stochastic conversion]
\label{cor:stochastic-conversion}
Let \(\lvert\psi\rangle,\lvert\phi\rangle\) be stabilizer
\(\AME(2m,q)\) states with \(m\geq2\), and let \(A_1,\ldots,A_{2m}\) be
linear operators on \(\mathbb C^q\) with
\[
 (A_1\otimes\cdots\otimes A_{2m})\lvert\psi\rangle=c\lvert\phi\rangle,
 \qquad c\neq0 .
\]
Then every \(A_i\) is a nonzero scalar multiple of a Clifford unitary.
Consequently equivalence under stochastic local operations and classical
communication, LU equivalence, and LC equivalence coincide for these
states; an exact single-copy conversion succeeds with probability one when
the states are LC equivalent and with probability zero otherwise; and
Corollary~\ref{cor:finite-recognition} decides which.
\end{corollary}

\begin{proof}
The one-party reductions of \(\phi\) have full rank, so each \(A_i\) is
invertible.  Write \(A_i=c_iU_ie^{H_i}\) with \(c_i\neq0\), \(U_i\)
unitary, and \(H_i\) Hermitian and traceless: this is the polar
decomposition, with the scalar chosen to remove the trace of the logarithm
of the positive factor.  Let \(H=\sum_iH_i\) act on party \(i\) through
\(H_i\), and put
\[
 f(t)=\log\bigl\|e^{tH}\lvert\psi\rangle\bigr\|^2,\qquad
 \lvert\psi_t\rangle
 =e^{tH}\lvert\psi\rangle\big/\bigl\|e^{tH}\lvert\psi\rangle\bigr\| .
\]
Then \(f'(t)=2\langle\psi_t\rvert H\lvert\psi_t\rangle\) and
\(f''(t)=4\operatorname{Var}_{\psi_t}(H)\geq0\).  The states
\(\psi_0=\psi\) and \(\psi_1\propto(\bigotimes_iU_i^\dagger)\phi\) both
have one-party reductions \(q^{-1}I\), and each \(H_i\) is traceless, so
\(f'(0)=f'(1)=0\).  Convexity forces \(f'\equiv0\) on \([0,1]\), hence
\(\operatorname{Var}_\psi(H)=0\) and \(H\lvert\psi\rangle=0\).  Expanding
\(\langle\psi\rvert H^2\lvert\psi\rangle=0\) with the two-party
reductions \(q^{-2}I\) of \(\psi\) gives
\[
 0=\sum_i\langle\psi\rvert H_i^2\lvert\psi\rangle
  +\sum_{i\neq j}\langle\psi\rvert H_iH_j\lvert\psi\rangle
  =q^{-1}\sum_i\Tr(H_i^2),
\]
so every \(H_i=0\) and \(A_i=c_iU_i\).
Theorem~\ref{thm:lu-lc-rigidity} makes each \(U_i\) Clifford.  A
successful branch of a stochastic protocol is a product operator as
displayed, so LC-inequivalent states cannot be converted with positive
probability, while LC-equivalent states are converted deterministically.
\end{proof}

Equality of the SLOCC and LU classes holds for every state with maximally
mixed one-party reductions by the Kempf--Ness theorem
\cite{KempfNess1979,GourWallach2011}, and for AME states explicitly
\cite[Corollary~1]{BurchardtRaissi2020}; the conversion dichotomy is
recorded in \cite{RajchelMieldziocEtAl2026}.  Kempf--Ness alone does not
make each factor unitary: the stabilizer of such a state in the product of
special linear groups is reductive and may contain a torus, as it does for
the Greenberger--Horne--Zeilinger state, which is one-uniform.  The
variance argument uses two-uniformity and applies verbatim to every
two-uniform state; the Clifford conclusion is the stabilizer-AME rigidity.
For qubit stabilizer states, compare the local-symmetry classification of
\cite{EnglbrechtKraus2020}.

\subsection{Encoder consequences}

\begin{proof}[Proof of Corollary~\ref{cor:transversal-clifford}]
The input marginal of each AME Choi state is \(q^{-1}I\), so one-leg
reshaping gives an isometry.  The projection of its stabilizer label group
onto the input Weyl-label space is surjective: otherwise a nonzero local
label orthogonal to the image would lie in the maximal isotropic space
\(L^\perp=L\), contradicting the absence of a one-party stabilizer.  The
kernel therefore has \(q^{2m-2}\) elements and stabilizes the image, making
it a one-logical-qudit stabilizer code.  AME corrects erasure of any
\(m-1\) outputs, so its distance is at least \(m\); the quantum Singleton
bound gives equality.  This is the standard perfect-tensor/quantum-MDS
correspondence used in \cite[Theorems~3.5 and 3.6]{Tan2026}.

Let \(\lvert\Phi\rangle=q^{-1/2}\sum_x\lvert x,x\rangle\), and use the
Choi convention
\(\lvert\psi\rangle=(I\otimes V_\psi)\lvert\Phi\rangle\) and
\(\lvert\phi\rangle=(I\otimes V_\phi)\lvert\Phi\rangle\).
Writing \(U_{\mathrm{phys}}=\bigotimes_iU_i\), the relation
\(U_{\mathrm{phys}}V_\psi=V_\phi L\) gives
\[
 (I\otimes U_{\mathrm{phys}})\lvert\psi\rangle
   =(L^T\otimes I)\lvert\phi\rangle .
\]
Hence
\[
 \bigl((L^T)^{-1}\otimes U_{\mathrm{phys}}\bigr)
   \lvert\psi\rangle=\lvert\phi\rangle .
\]
Theorem~\ref{thm:lu-lc-rigidity} makes every displayed local factor
Clifford.  Entrywise conjugation preserves the finite-field Weyl group,
since \(\overline{X(a)}=X(a)\) and
\(\overline{Z(b)}=Z(-b)\); it therefore preserves the Clifford
normalizer.  Taking adjoints preserves the normalizer as well, and
\(U^T=(\overline U)^\dagger\).  Thus transposition preserves the
finite-field Clifford group, and \(L\) is Clifford.
\end{proof}

The recognition procedure is constructive for these encoders.  Given check
matrices for the two normalized Choi states,
Corollary~\ref{cor:finite-recognition} decides whether a transversal conversion
exists and returns compact symplectic-frame witnesses within the same
arithmetic bound.  If full stabilizer tableaux and chosen Clifford lifts are
also supplied, one further linear solve returns a compact Clifford--Pauli
conversion: the reference frame \(R\) gives logical frame
\((R^{\mathsf T})^{-1}\), and
Lemma~\ref{lem:pauli-phase-correction} allows the character repair to be
supported on any prescribed \(m\) physical outputs.  Thus the repair need not
alter the logical factor.  Moreover, \(m\) is the smallest size of a
prescribed carrier that works uniformly for every discrepancy.  This bound
does not include expansion into dense unitary matrices.
